% =====================================================================
%  5.1 Sinus, cosinus og tangens — fasit
% =====================================================================

\fk{Del 1 — Sinus, cosinus og tangens i rettvinklede trekanter}
\fo{1.1} $\sin v = \frac{5}{13}$ \ $\cos v = \frac{12}{13}$ \ $\tan v = \frac{5}{12}$
\fo{1.2} a) $0,8 \cdot 15 = 12$ \ b) $\sqrt{15^2 - 12^2} = 9$ \ c) $\cos v = \frac{9}{15} = \frac{3}{5}$
\fo{1.3} a) $0,423$ \ b) $0,438$ \ c) $1,428$ \ d) $63,4^\circ$ \ e) $\sin^{-1}$ er den omvendte funksjonen som gir vinkelen, ikke $\frac{1}{\sin}$. Riktig: $\sin^{-1}(0,4) \approx 23,6^\circ$.
\fo{1.4} a) Hypotenus: $\sqrt{9^2 + 40^2} = 41$. \ b) $\sin v = \frac{9}{41}$, $\cos v = \frac{40}{41}$, $\tan v = \frac{9}{40}$
\fo{1.5} a) $12,0 \cdot \cos 37^\circ \approx 9,6$ cm \ b) $9,0 \cdot \tan 42^\circ \approx 8,1$ cm \ c) $\frac{5,5}{\sin 33^\circ} \approx 10,1$ cm
\fo{1.6} a) $\cos v = \frac{1,5}{5,0}$, så $v \approx 72,5^\circ$. \ b) $\sqrt{5,0^2 - 1,5^2} \approx 4,8$ m. Kontroll: $5,0 \cdot \sin 72,5^\circ \approx 4,8$ m.
\fo{1.7} Den første feilen er overgangen til $x = 14 \cdot (-0,667)$. Kalkulatoren stod i RAD, og $\cos 40$ i radianer er $-0,667$.
Riktig i DEG: $\cos 40^\circ \approx 0,766$, så $x = 14 \cdot 0,766 \approx 10,7$. En lengde kan aldri være negativ.
\fo{1.8} a) Alltid. En katet er kortere enn hypotenusen, så $\frac{\text{motstående}}{\text{hypotenus}} < 1$.
b) Noen ganger. $\tan v < 1$ når motstående katet er kortere enn hosliggende. Katetene $3$ og $4$ gir $\frac34 < 1$, men byttes de, blir $\frac43 > 1$.
c) Aldri. Med sidene $a$, $b$ og $c$ er $\sin v + \cos v = \frac{a+b}{c}$. To sider i en trekant er til sammen lengre enn den tredje, så $a + b > c$ og summen er større enn 1.
\fo{1.9} Med motstående $a$, hosliggende $b$ og hypotenus $c$: $\frac{\sin v}{\cos v} = \frac{a/c}{b/c}$. Vi multipliserer med $c$ i teller og nevner, og får $\frac{a}{b}$. Det er $\tan v$.
\fo{1.10} Katetene er $14$ og $48$. Trekanten $7$, $24$, $25$ har $\tan v = \frac{7}{24}$ og hypotenus $25$. Vi dobler alle sidene og får $14$, $48$, $50$. Kontroll: $14^2 + 48^2 = 196 + 2304 = 2500 = 50^2$.
Det finnes bare én slik trekant, fordi to rettvinklede trekanter med samme vinkel er formlike, og når hypotenusen er gitt, er målestokken gitt.
\fo{1.11} a) $\sin v = \frac{6}{10} = \frac35$ og $\sin v = \frac{9}{15} = \frac35$. Like svar.
b) Trekantene har vinklene $v$, $90^\circ$ og $90^\circ - v$, så de er formlike. Forholdet mellom to sider endres ikke når alle sidene ganges med samme tall. Derfor avhenger $\sin v$ bare av $v$.
\fo{Sjekk 1} a) $\cos v = \frac{24}{25}$ (hosliggende er $\sqrt{25^2 - 7^2} = 24$) \ b) $0,77$ \ c) Trekantene er formlike fordi vinklene er like. Da er forholdet mellom to sider det samme i begge trekantene.

\fk{Del 2 — Eksakte verdier og enhetssirkelen}
\fo{2.1} a) $\frac12$ \ b) $\frac12$ \ c) $1$ \ d) $\frac{\sqrt2}{2}$
\fo{2.2} Motstående: $14 \cdot \sin 30^\circ = 7$. Hosliggende: $14 \cdot \cos 30^\circ = 7\sqrt3$. Kontroll: $7^2 + (7\sqrt3)^2 = 49 + 147 = 196 = 14^2$.
\fo{2.3} a) $1$ \ b) $1$ \ c) $-1$ \ d) $0$
\fo{2.4} a) $\frac{\sqrt3}{2}$ \ b) $\frac{\sqrt2}{2}$ \ c) $-\frac{\sqrt2}{2}$ \ d) $-\frac12$
\fo{2.5} a) $v \approx 53,1^\circ$ og $v \approx 126,9^\circ$ \ b) $v \approx 104,5^\circ$ (bare én) \ c) $v \approx 66,9^\circ$ og $v \approx 113,1^\circ$
Kalkulatoren gir bare den spisse løsningen for sinus. Den andre er $180^\circ$ minus den første.
\fo{2.6} a) $\sin 90^\circ = 1$, men $2\sin 45^\circ = 2 \cdot \frac{\sqrt2}{2} = \sqrt2$. Påstanden stemmer ikke generelt. \ b) $v = 0^\circ$: $\sin 0^\circ = 0 = 2 \cdot 0$.
\fo{2.7} Den første feilen er overgangen til $\cos v = \frac{15}{17}$. Siden $v$ er stump, er cosinus negativ. Riktig er $\cos v = -\frac{15}{17}$.
Da er $v = \cos^{-1}\!\left(-\frac{15}{17}\right) \approx 151,9^\circ$. Kontroll: $\sin 151,9^\circ \approx 0,471 \approx \frac{8}{17}$.
\fo{2.8} a) Noen ganger. Positiv for $v < 90^\circ$ (punktet har positiv $x$-koordinat), men $\cos 120^\circ$ er negativ.
b) Noen ganger. Likheten betyr at $x$- og $y$-koordinaten er like. Det skjer bare for $v = 45^\circ$, der begge er $\frac{\sqrt2}{2}$. For $v > 90^\circ$ er $\cos v < 0 \le \sin v$.
c) Aldri. Cosinus er $x$-koordinaten. Samme $x$-koordinat på den øvre halvsirkelen gir samme punkt, og dermed samme vinkel. Altså er $u = v$.
\fo{2.9} Enten er $v_1 = v_2$, eller så er $v_1 = 180^\circ - v_2$. Det siste er umulig i en trekant, for da er $v_1 + v_2 = 180^\circ$, og den tredje vinkelen ville vært $0^\circ$. Altså er $v_1 = v_2$, og trekanten er likebeint.
\fo{Sjekk 2} a) $\sqrt3$ \ b) $v = 60^\circ$ og $v = 120^\circ$ \ c) $\sin v = \sqrt{1 - (-0,6)^2} = 0,8$. Punktet ligger på enhetssirkelen, så $x^2 + y^2 = 1$, og $y > 0$.

\fk{Ekstra utfordring}
\fo{U1} Rampelengde: $\frac{1,5}{\sin 3^\circ} \approx 28,7$ m. Langs bakken: $\frac{1,5}{\tan 3^\circ} \approx 28,6$ m.
\fo{U2} a) $\frac34 - \frac14 = \frac12$ \ b) $2 \cdot \frac{\sqrt2}{2} \cdot \frac{\sqrt2}{2} = 1$ \ c) $\frac{\sqrt2}{2} \cdot \left(-\frac{\sqrt2}{2}\right) = -\frac12$
\fo{U3} a) $h = 8 \cdot \sin 60^\circ = 4\sqrt3$. Areal: $\frac12 \cdot 8 \cdot 4\sqrt3 = 16\sqrt3$.
b) Høyden er $s \cdot \sin 60^\circ = \frac{\sqrt3}{2}s$. Areal: $\frac12 \cdot s \cdot \frac{\sqrt3}{2}s = \frac{\sqrt3}{4}s^2$.
\fo{U4} a) Funksjonene i \lstinline[style=laereverk,numbers=none,frame=none]|math| regner i radianer. \lstinline[style=laereverk,numbers=none,frame=none]|math.radians| gjør grader om til radianer. Uten den gir for eksempel \lstinline[style=laereverk,numbers=none,frame=none]|math.sin(30)| svaret $-0,988$.
b) Parene er $(0,180)$, $(15,165)$, $(30,150)$, $(45,135)$, $(60,120)$ og $(75,105)$. Hvert par summerer til $180^\circ$. Punktene på enhetssirkelen ligger speilvendt om $y$-aksen, og har samme $y$-koordinat.
c) Flyttall er bare tilnærminger. \lstinline[style=laereverk,numbers=none,frame=none]|math.radians(180)| er et avrundet $\pi$, så sinus blir en svært liten verdi, omtrent $1,2 \cdot 10^{-16}$. \lstinline[style=laereverk,numbers=none,frame=none]|round| runder den til $0$.
\fo{U5} La $x$ være avstanden fra første punkt til treet og $h$ høyden. Da er $\tan 50^\circ = \frac{h}{x}$ og $\tan 30^\circ = \frac{h}{x+20}$.
Fra første likning er $x = \frac{h}{\tan 50^\circ}$. Setter vi inn i den andre, får vi $h\left(\frac{1}{\tan 30^\circ} - \frac{1}{\tan 50^\circ}\right) = 20$. Da er $h \approx 22,4$ m.
\fo{U6} a) Punktet $(\cos v,\ \sin v)$ ligger på sirkelen med radius 1. Pytagoras for avstanden til origo gir $\cos^2 v + \sin^2 v = 1^2$.
b) Sett $u = 180^\circ - v$. Da er $\sin v = \sin u$ og $\cos v = -\cos u$. Kvadrerer vi, blir $\sin^2 v + \cos^2 v = \sin^2 u + \cos^2 u = 1$, siden $u$ er spiss.
\fo{U7} Fra 1.9 er $\tan v = \frac{\sin v}{\cos v}$. For spisse $v$ er $0 < \cos v < 1$, så $\frac{\sin v}{\cos v} > \sin v$. Dermed er $\sin v < \tan v$.
Eksempel på en annen påstand: $\cos v < \tan v$. Den er noen ganger sann: for $v = 30^\circ$ er $\cos v \approx 0,87 > \tan v \approx 0,58$, men for $v = 60^\circ$ er $\cos v = 0,5 < \tan v \approx 1,73$. Flere påstander er riktige.
\fo{U8} Noen ganger. Likningen betyr at $x + y = 1$ for punktet $(x,\ y)$ på enhetssirkelen. Linjen $x + y = 1$ skjærer sirkelen i $(1,\ 0)$ og $(0,\ 1)$. Det gir $v = 0^\circ$ og $v = 90^\circ$. For $v = 180^\circ$ er summen $0 + (-1) = -1$.
